\documentclass[12pt]{article}

% Language setting
\usepackage[ngerman]{babel}

% Set page size and margins
\usepackage[a4paper,top=4cm,bottom=2cm,left=2.5cm,right=2.5cm,marginparwidth=1.75cm]{geometry}
\usepackage{float}

% Useful packages
\usepackage{amsmath}
\usepackage{graphicx}
\usepackage[colorlinks=true, allcolors=blue]{hyperref}
\usepackage[figure]{hypcap}
\usepackage[inkscapeformat=png]{svg}
\usepackage{fancyhdr}      
\usepackage{geometry}
\usepackage{cite}
\usepackage{wrapfig}
\usepackage{lipsum}


\setcounter{secnumdepth}{4}

% Page layout
\geometry{margin=1in}
\pagestyle{fancy}
\setlength{\headheight}{50.5120pt}  
\setlength{\headsep}{20pt}     % Abstand zwischen Header und Text


% Header setup
\fancyhf{} % header und footer aufräumen

%\fancyhead[L]{\includegraphics[width=2.5cm,height=2.5cm,keepaspectratio]{}}
\fancyhead[R]{
  \begin{tabular}{r}
    \textbf{Name:} Stolz, Katzenski, Behm \\
    \textbf{Abgabe:} 21.Juni 2026 
  \end{tabular}
}
 %\addtolength{\topmargin}{-20.51192pt}
% Footer (optional)
\fancyfoot[C]{\thepage}


\begin{document}

\vspace*{0.01cm}

\begin{center}
    \Large
    \textbf{Projekt Dokumentation: Roguelite CLI  }
        
    \vspace{0.4cm}
    \large
   
       
    %\vspace{0.9cm}
    %\textbf{Abstract}
\end{center}


\tableofcontents 
\vspace*{0.5cm}
\noindent 

% -------------------------------------- %

\section{Einleitung}

% -------------------------------------- %

\section{Planung}

% -------------------------------------- %

\subsection{Kanban Board}


% -------------------------------------- %

\subsection{MoScOw Methode}


% -------------------------------------- %

\subsection{In-Plane Switching}



% -------------------------------------- %

\section{Spiel Anleitung}


% -------------------------------------- %


\end{document}


% -------------------------------------- nützliches Zeug

%\section
%\subscection
%\subsubsection

%\begin{figure}[H] %[H] 
%\centering
%\includegraphics[width=0.3\textwidth]{ZoneBitRecording.png}
%\caption{\label{fig:ZBR}Sektoren Einteilung einer HDD mit Zone-Bit-Recording}
%\end{figure}

%\begin{wrapfigure}{r}{5.5cm}
%\includegraphics[width=5.5cm]{sandwich.png}
%\caption{Sudo make a sandwich}\label{wrap-fig:1}
%\end{wrapfigure} 

%\begin{figure}[H]
%\centering
%\begin{minipage}{.5\textwidth}
%  \centering
%  \includegraphics[width=.8\linewidth]{LCD-Hintergrundbeleuchtung_mit_Edge-LEDs.JPG}
%  \caption{LCD mit Edge LED \cite{wikipediaHintergrundbeleuchtung}}
%  \label{fig:edgeLED}
%\end{minipage}%
%\begin{minipage}{.5\textwidth}
%  \centering
%  \includegraphics[width=.8\linewidth]{LCD-Hintergrundbeleuchtung_mit_LEDs.JPG}
%  \caption{Full Array LED \cite{wikipediaHintergrundbeleuchtung}}
%  \label{fig:fullArrayLED}
%\end{minipage}
%\end{figure}


%\begin{enumerate}
%    \item 
%\end{enumerate}

%\noindent
%\space
%\\  -> Absatz
%\newline
% \glqq \grqq -> Deutsche Anführungszeichen
%\cite{x, y}


%\nocite{*}  %-> alle Quellen übernehmen ins Inhaltsverzeichnis, auch wenn nicht direkt zitiert
%\bibliographystyle{alpha}
%\bibliography{sample}

